\section{Synopsis}
\label{sec:synopsis}

A nested non-stationary dynamic Bayesian network is built from a ladder of
temporal resolutions, two maps between adjacent rungs, one latent state and one
lagged graph per rung. The densities, the estimator and the proofs follow in
\cref{sec:model,sec:learning,sec:identifiability}.

\paragraph{The ladder.}
An event log is \(\mathcal D=\{(\ell_i,s_i,d_i)\}\): a label, a start, a
duration. Fix nested partitions of the study horizon indexed by
\(k=1,\dots,K\), finest to coarsest, and read off each epoch the number of
episodes that start in it and the duration they carry,
\begin{align}
 Z^{(k)}=\bigl(X^{(k)},Y^{(k)}\bigr),
 \qquad
 X_t^{(k,\ell)}=\#\{i:\ell_i=\ell,\ s_i\in I_t^{(k)}\},
 \qquad
 Y_t^{(k,\ell)}=\!\!\sum_{i:\,s_i\in I_t^{(k)}}\!\! d_i .
 \label{eq:syn_project}
\end{align}
One log, \(K\) views of it. Write \(\child_k(c)\) for the children of a
coarse epoch \(c\).

\paragraph{Folding and unfolding.}
Two maps move between adjacent views. The \emph{fold} is deterministic and
known; the \emph{unfold} is the model,
\begin{align}
 \underbrace{\bigl(\mathcal A_kz\bigr)_c=\!\!\sum_{t\in\child_k(c)}\!\! z_t}_{\text{fold, fine}\to\text{coarse}},
 \qquad
 \underbrace{Z^{(k)}\sim W_k\bigl(\cdot\mid Z^{(k+1)},U^{(k)},C^{(k)}\bigr)}_{\text{unfold, coarse}\to\text{fine}},
 \qquad
 \operatorname{supp}W_k=\mathcal A_k^{-1}\bigl(Z^{(k+1)}\bigr).
 \label{eq:syn_foldunfold}
\end{align}
The unfold is a distribution over the \emph{fiber}: over all the ways a coarse
total can be dealt out among its children. Generation starts at the top and
walks down,
\begin{align}
 Z^{(K)}\sim E_K\bigl(\cdot\mid U^{(K)},C^{(K)}\bigr),
 \qquad
 Z^{(K-1)}\sim W_{K-1},\ \dots,\ Z^{(1)}\sim W_1 ,
 \label{eq:syn_generate}
\end{align}
where \(U^{(k)}=(S^{(k)},H^{(k)})\) is a per-scale latent state, a discrete
regimen for persistent drift and a continuous factor for same-epoch
co-movement, and \(C^{(k)}\) holds calendar quantities as exogenous nodes.
Because every unfold lands inside the fiber it came from,
\begin{align}
 \mathcal A_k\circ W_k=\mathrm{id}
 \qquad\text{almost surely, hence}\qquad
 \mathcal A_{j-1}\cdots\mathcal A_kZ^{(k)}=Z^{(j)}\ \ \text{for all }k<j .
 \label{eq:syn_coherence}
\end{align}
Every draw, every forecast and every simulation reconciles exactly at every
resolution. Nothing is projected afterwards.

\paragraph{Shadow and detail.}
Each scale carries a lagged graph \(\mathcal G^{(k)}\) and a predictor
\(\eta^{(k)}\) built from lagged parents, latents and calendar. Split that
predictor along the blocks, with \(\Pi_k\) the block average,
\begin{align}
 \eta^{(k)}
 =\underbrace{\Pi_k\eta^{(k)}}_{\text{shadow}}
 \;\oplus\;
 \underbrace{(I-\Pi_k)\eta^{(k)}}_{\text{detail}} .
 \label{eq:syn_split}
\end{align}
The unfold is blind to its own shadow: adding a constant to every child of one
block leaves \(W_k\) unchanged, because the split it induces is unchanged.
Therefore
\begin{center}
\emph{the detail at scale \(k\) is identified at scale \(k\); the shadow at
scale \(k\) is identified only above it.}
\end{center}
This single sentence organizes the rest, and it has a spectral reading: folding
is a low-pass filter, so the resolutions partition frequency and each one owns
its own band. Periodic structure is then carried by a few Fourier harmonics,
each landing in the band that contains its frequency, and localized structure by
Haar coefficients on the block tree, where a level shift is simply a large one.
A changepoint therefore has a home scale exactly as a calendar quantity does.
Every parameter has a resolution that can see it. A latent regimen that does not switch inside a block is pure shadow
there, so below the coarsest scale a regimen must act on a feature that varies
within the block. A calendar quantity is pure shadow above the scale at which
its period fits, which is what gives each one a home: hour of day at the hour,
weekday at the day, a study trend at the week. And two adjacent scales share
exactly their shadow, which is the only thing a cross-resolution prior may tie.

\paragraph{The nested walk.}
Learning is one walk up the ladder and one walk down, repeated. The ascent
carries evidence, the descent carries context,
\begin{align}
 \text{ascent }(1\!\to\!K):\ \text{fold the details into sufficient statistics};
 \qquad
 \text{descent }(K\!\to\!1):\ \text{pass shadows down as priors} .
 \label{eq:syn_walk}
\end{align}
Both directions are necessary, and \cref{eq:syn_split} says why: neither scale
can identify what the other owns. One sweep updates the latent posterior
\(q(U)\), then the parameters, then proposes graph moves, each step
non-decreasing on one penalized bound
\begin{align}
 \mathcal J
 =\mathbb E_q\bigl[\log p(Z,U,\mathcal G,\Theta\mid C)\bigr]
 +\mathbb H[q]
 -\sum_k\Bigl[\lambda_k^{\mathrm{rel}}\!\!\sum_{r}z_{r,k}
 +\lambda_k^{\mathrm{rep}}\!\!\sum_{r}z_{r,k}\bigl(|\mathcal R_k(r)|-1\bigr)\Bigr] .
 \label{eq:syn_objective}
\end{align}
A scientific relation \(r\) is bought once at price
\(\lambda_k^{\mathrm{rel}}\); its representations, the choice of lag, quantity
and basis, are cheap refinements inside it. Pricing each representation
separately is what makes a simple relation lose to its own alternatives.

\paragraph{The cascade.}
The descent carries one more thing, and it is what makes the search
affordable. A relation whose lag is shorter than a coarse epoch cannot appear
there as a lag; it appears as same-epoch co-movement. Coarse co-movement is
therefore the shadow of a fine lagged edge, and it orders where to look,
\begin{align}
 \rho_{a\to b}^{(k)}
 =\max_{j>k}\ \bigl|\operatorname{cor}\bigl(Z^{(j,a)},Z^{(j,b)}\bigr)\bigr|
 \quad\longrightarrow\quad
 \text{visit candidates at scale }k\text{ in decreasing }\rho .
 \label{eq:syn_cascade}
\end{align}
This is an ordering and never a filter. Nothing is removed, so at full depth
the accepted set is exactly the unscreened one and the traversal depth is a
speed control with no effect on what is reachable. Coarse regimens order the
search the same way, by naming the epochs in which a relation is worth
examining at all.

\paragraph{What comes back.}
A fit returns, at every resolution, a graph, its conditional distributions, and
the latent regimen and factor paths. Running \cref{eq:syn_generate} forward
from a filtered state produces a forecast; running it without conditioning
produces a synthetic event log. Both are coherent by \cref{eq:syn_coherence},
so a claim made at one resolution is automatically a claim at all of them, and
a synthetic log can be re-projected through \cref{eq:syn_project} and compared
with the real one scale by scale. \Cref{sec:results:arm1} runs that comparison
on an observed log, and the pattern it finds is the one this synopsis predicts:
what a fit reproduces is coarse-scale drift, and where the two logs separate is
inside the day.
